% 原书第二篇2.7节240--245，用户PDF页99--100，印刷页90。
% 240--242使用经典硬球碰撞近似，超稀薄极限单独列模型。
\thermotransportproblem{240}{90}
已知标准状况下氦气的动力黏度系数，
利用分子运动论计算氦原子的有效直径。
\exsolution
在经典硬球近似下，
$\lambda=(\sqrt2\pi d^2n_d)^{-1}$，
而动理论的黏度公式是
$\eta=\rho\bar v\lambda/3$。
令单原子质量为
$m=M/N_A$，
$\rho=n_dm$，
得到
\[
\eta=\frac{n_dm\bar v}
{3\sqrt2\pi d^2 n_d}
=\frac{m\bar v}{3\sqrt2\pi d^2}.
\]
因此
\[
\boxed{d=\sqrt{
\frac{m\bar v}{3\sqrt2\pi\eta}
}},\qquad
\bar v=\sqrt{\frac{8RT_0}{\pi M}}.
\]
标准状况下取氦气黏度
$\eta\simeq(1.8$--$2.0)\times10^{-5}\,
\mathrm{Pa\,s}$、
$M=0.0040\,\mathrm{kg/mol}$，
得到
\[
\boxed{d\simeq1.8\times10^{-10}\,\mathrm m
=0.18\,\mathrm{nm}}.
\]
\excheck 原题没有明列黏度的数值，
需要使用标准状况氦气物性值；
本结果属于所采用硬球碰撞模型的\emph{有效直径}，
不应与量子化学意义上的原子硬边界完全等同。

\thermotransportproblem{241}{90}
标准状况下，
氦气的热导率是氩气热导率的
$8.7$ 倍。
求氩原子与氦原子有效直径的比
$d_{\rm Ar}/d_{\rm He}$。
\exsolution
气体动理论给
$\kappa=(1/3)\rho c_V\bar v\lambda$。
同为单原子气体，
质量比热
$c_V=(3/2)k_B/m$，
而
$\rho=n_dm$、
$\lambda=(\sqrt2\pi d^2n_d)^{-1}$，
所以
\[
\kappa\propto
(n_dm)\frac1m
\sqrt{\frac{T}{m}}\,
\frac1{d^2n_d}
\propto\frac{\sqrt T}{d^2\sqrt m}.
\]
在相同标准温度下，
\[
\frac{\kappa_{\rm He}}{\kappa_{\rm Ar}}
=\frac{d_{\rm Ar}^2\sqrt{m_{\rm Ar}}}
{d_{\rm He}^2\sqrt{m_{\rm He}}}.
\]
解出
\[
\boxed{\frac{d_{\rm Ar}}{d_{\rm He}}
=\sqrt{
\frac{\kappa_{\rm He}}{\kappa_{\rm Ar}}
\sqrt{\frac{m_{\rm He}}{m_{\rm Ar}}}}}.
\]
用 $M_{\rm He}\simeq4\,\mathrm{g/mol}$、
$M_{\rm Ar}\simeq40\,\mathrm{g/mol}$：
\[
\boxed{d_{\rm Ar}/d_{\rm He}\simeq1.66
\approx1.7}.
\]
\excheck 两气体都是单原子气体，
故其热容量在比例中相消；
不可漏掉原子质量的平方根因子。

\thermotransportproblem{242}{90}
标准状况下，
两根同轴长圆筒之间充满氦气。
圆筒平均半径为 $R$，
间隙 $\Delta R\ll R$，
外圆筒以小角速度 $\omega$ 转动，
内筒静止。
求内筒每单位长度所受摩擦力矩。
当 $\Delta R=6.0\,\mathrm{mm}$ 时，
为使力矩减小 $n=10$ 倍，
应该将氦气压强降到什么量级？
\exsolution
在连续层流黏滞近似下，
两壁相对切向速度为
$\Delta u\simeq\omega R$，
薄隙剪切应力
$\tau_w\simeq\eta\Delta u/\Delta R
=\eta\omega R/\Delta R$。
内筒每单位长度侧面积为 $2\pi R$，
故摩擦力矩每单位长度
\[
\boxed{\mathcal N_0\simeq
\tau_w(2\pi R)R
=\frac{2\pi\eta\omega R^3}{\Delta R}}.
\]
大幅降低气压至自由程超过板间隙，
进入超稀薄气体模型，题给
剪切应力
$\tau_{\rm rare}=\rho\bar v\,\Delta u/6$，
因此
\[
\mathcal N_{\rm rare}
=\frac{\pi}{3}\rho\bar v\,\omega R^3.
\]
要求
$\mathcal N_{\rm rare}=\mathcal N_0/n$：
\[
\rho\bar v=\frac{6\eta}{n\Delta R}.
\]
用
$\eta\simeq\rho\bar v\lambda/3$、
$\lambda=k_BT/(\sqrt2\pi d^2p)$
及 $\rho\bar v\propto p$，
消去密度得到压强的量级：
\[
\boxed{p\simeq\frac{\sqrt2\,k_BT}
{\pi d^2n\Delta R}}.
\]
若氦原子有效直径取
$d\simeq0.18\,\mathrm{nm}$、
$T\simeq273\,\mathrm K$、$n=10$，
\[
\boxed{p\ \text{为 }1\,\mathrm{Pa}\text{ 左右}
\ (\text{原书近似值 }0.7\,\mathrm{Pa})}.
\]
\excheck 连续黏性与自由分子流分别适用于
$\lambda\ll\Delta R$、
$\lambda\gg\Delta R$；
在过渡 Knudsen 数区间需要更精细的动力学模型，
这里的压强只作为量级估计。

\thermotransportproblem{243}{90}
一个气体充满同轴长圆筒之间的空间，
内外半径为 $R_1<R_2$。
外筒以角速度 $\omega$ 缓慢转动，内筒静止。
已知内筒每单位轴向长度上的
黏滞摩擦力矩大小为 $\mathcal N_1$。
设半径 $r$ 处剪切应力
$\tau(r)=\eta r\,\mathrm d\Omega(r)/\mathrm dr$，
求气体动力黏度 $\eta$。
\exsolution
稳定轴对称旋转层流中，
通过每个同轴圆柱面的力矩每单位长度
必须相同：
\[
\mathcal N_1=2\pi r^2\tau(r)
=2\pi\eta r^3
\frac{\mathrm d\Omega}{\mathrm dr}.
\]
因此
$\mathrm d\Omega/\mathrm dr=
\mathcal N_1/(2\pi\eta r^3)$。
由边界
$\Omega(R_1)=0$、
$\Omega(R_2)=\omega$：
\[
\omega=\int_{R_1}^{R_2}
\frac{\mathcal N_1}{2\pi\eta r^3}\,\mathrm dr
=\frac{\mathcal N_1}{4\pi\eta}
\left(\frac1{R_1^2}-\frac1{R_2^2}\right).
\]
解得
\[
\boxed{\eta=\frac{\mathcal N_1}{4\pi\omega}
\left(\frac1{R_1^2}-\frac1{R_2^2}\right)}.
\]
\excheck 必须保留半径平方的倒数差，
且 $\mathcal N_1$ 是\emph{单位长度力矩}，
对应量纲为 $\mathrm{N}$。

\thermotransportproblem{244}{90}
两块半径均为 $a$ 的圆盘具有公共轴线，
间距 $h\ll a$，
其中一块静止，
另一块以小角速度 $\omega$ 旋转。
盘间气体黏度为 $\eta$。
求静止圆盘所受黏滞摩擦力矩。
\exsolution
在距中心半径 $r$ 处，
旋转盘的切向速率为 $\omega r$。
窄间隙中的线性速度梯度近似为
$\partial u_\theta/\partial z=\omega r/h$，
剪切应力
\[
\tau(r)=\eta\frac{\omega r}{h}.
\]
静止圆盘上的环带面积
$\mathrm dA=2\pi r\,\mathrm dr$，
相应力矩微元
\[
\mathrm d\mathcal N=
r\tau(r)\,\mathrm dA
=\frac{2\pi\eta\omega}{h}r^3\,\mathrm dr.
\]
积分 $0\le r\le a$：
\[
\boxed{\mathcal N=
\frac{2\pi\eta\omega}{h}
\int_0^a r^3\,\mathrm dr
=\frac{\pi\eta\omega a^4}{2h}}.
\]
\excheck 盘面只有\emph{一个面对间隙中的气体}
承受该剪切力矩，
不能像整个转盘双侧浸液那样再乘2。

\thermotransportproblem{245}{90}
仍是上一题的两块平行圆盘，
但盘间改为超稀薄气体，
摩尔质量为 $M$、
温度 $T$、压强 $p$。
求静止圆盘所受摩擦力矩。
\exsolution
超稀薄气体中自由程远大于板间距，
不再适用连续剪切
$\eta\omega r/h$。
按本节给出的自由分子流近似，
两板相对切向速率为
$\Delta u=\omega r$，
气体对单位面积的剪切应力为
\[
\tau(r)=\frac16\rho\bar v\,\omega r.
\]
半径为 $r$ 的环带面积
$2\pi r\,\mathrm dr$，故
\[
\mathrm d\mathcal N
=r\tau(r)(2\pi r\,\mathrm dr)
=\frac{\pi}3\rho\bar v\omega r^3\,\mathrm dr.
\]
积分：
\[
\mathcal N=\frac{\pi}{12}\rho\bar v\omega a^4.
\]
由理想气体关系
$\rho=pM/(RT)$，
以及
$\bar v=\sqrt{8RT/(\pi M)}$，
得到
\[
\boxed{\mathcal N=
\frac{\omega a^4p}{3}
\sqrt{\frac{\pi M}{2RT}}}.
\]
\excheck 与连续层流结果不同，
本题的力矩不含板间距 $h$，
而随稀薄气体压强 $p$ 线性变化，
符合自由分子输运极限。
